设离散型随机变量 $X\in\{0,1,2,\dots\}$，则 $G_X(z)=\sum_{n\ge0}P(X=n)z^n$。

记 $[z^n]G_X(z)$ 为 $G_X(z)$ 中 $z^n$ 的系数，$G_X^{(k)}$ 为 $k$ 阶导数，$X^{\underline k}=X(X-1)\cdots(X-k+1)$（$k=0$ 时为 $1$）。$\begin{Bmatrix}k\\i\end{Bmatrix}$ 是第二类斯特林数。

\textbf{基本结论}：

$P(X=n)=[z^n]G_X(z)$

$E[X]=G_X'(1)$

$E[X^{\underline k}]=G_X^{(k)}(1)$

$E[X^k]=\sum_{i=0}^k\begin{Bmatrix}k\\i\end{Bmatrix}G_X^{(i)}(1)$

$\operatorname{Var}(X)=G_X''(1)+G_X'(1)-G_X'(1)^2$

\textbf{独立变量}：

若 $X_1,\dots,X_n$ 相互独立，且 $S=X_1+\cdots+X_n$，则 $G_S(z)=\prod_{i=1}^nG_{X_i}(z)$。

若各 $X_i$ 与 $X$ 同分布且相互独立，则 $G_S(z)=G_X(z)^n$。

\textbf{随机次数求和}：

若各 $X_i$ 与 $X$ 同分布且相互独立，$N$ 与整个序列 $(X_i)$ 独立，令 $S=X_1+\cdots+X_N$，则 $G_S(z)=G_N(G_X(z))$。

\textbf{常见离散分布的 PGF}：

以下约定 $q=1-p$，$k\ge0$。几何分布与负二项分布均基于成功概率为 $p$ 的独立重复试验：前者记录首次成功前的失败次数，后者记录第 $r$ 次成功前的失败次数。


1. 伯努利分布

$P(X=k)=p^k(1-p)^{1-k},\quad k\in\{0,1\}$

$G(z)=1-p+pz$

$G'(z)=p$

$G^{(k)}(z)=0,\quad k\ge 2$

2. 二项分布

$P(X=k)=\dbinom{n}{k}p^k(1-p)^{n-k},\quad k=0,1,\dots,n$

$G(z)=(1-p+pz)^n=\displaystyle\sum_{k=0}^{n}\binom{n}{k}p^k(1-p)^{n-k}z^k$

$G'(z)=np(1-p+pz)^{n-1}$

$G^{(k)}(z)=n^{\underline{k}}p^k(1-p+pz)^{n-k}$

3. 几何分布

$P(X=k)=(1-p)^kp,\quad k=0,1,2,\dots$

$G(z)=\dfrac{p}{1-(1-p)z}=\displaystyle\sum_{k=0}^{\infty}(1-p)^kp z^k$

$G'(z)=\dfrac{p(1-p)}{[1-(1-p)z]^2}$

$G^{(k)}(z)=\dfrac{k!\,p(1-p)^k}{[1-(1-p)z]^{k+1}}$

4. 负二项分布

$P(X=k)=\binom{k+r-1}{k}p^r(1-p)^k,\quad k=0,1,2,\dots$

$G(z)=\left(\dfrac{p}{1-(1-p)z}\right)^r=\displaystyle\sum_{k=0}^{\infty}\binom{k+r-1}{k}p^r(1-p)^k z^k$

$G'(z)=\dfrac{rp^r(1-p)}{[1-(1-p)z]^{r+1}}$

$G^{(k)}(z)=r^{\overline{k}}p^r(1-p)^k[1-(1-p)z]^{-r-k}$

5. 泊松分布

$P(X=k)=e^{-\lambda}\dfrac{\lambda^k}{k!},\quad k=0,1,2,\dots$

$G(z)=e^{\lambda(z-1)}=\displaystyle\sum_{k=0}^{\infty}e^{-\lambda}\dfrac{\lambda^k}{k!}z^k$

$G'(z)=\lambda e^{\lambda(z-1)}$

$G^{(k)}(z)=\lambda^k e^{\lambda(z-1)},\quad k\ge 0$

6. 离散均匀分布

$P(X=k)=\dfrac1{m+1},\quad k=0,1,\dots,m$

$G(z)=\dfrac{1-z^{m+1}}{(m+1)(1-z)}=\dfrac1{m+1}\displaystyle\sum_{j=0}^{m} z^j$

$G'(z)=\dfrac1{m+1}\displaystyle\sum_{j=1}^{m} jz^{j-1}$

$G^{(k)}(z)=\dfrac1{m+1}\displaystyle\sum_{j=k}^{m} j^{\underline{k}}z^{j-k}$

7. 超几何分布

总体有 $N$ 个物品，其中 $K$ 个为目标物品。从中不放回地等概率抽取 $n$ 个，$X$ 表示抽到的目标物品数。

设 $k_{\min}=\max(0,n-(N-K)),\qquad k_{\max}=\min(K,n)$

$P(X=k)=\dfrac{\binom{K}{k}\binom{N-K}{n-k}}{\binom{N}{n}},\quad k=k_{\min},k_{\min}+1,\dots,k_{\max}$

$G(z)=\dfrac1{\binom{N}{n}}\displaystyle\sum_{k=k_{\min}}^{k_{\max}}\binom{K}{k}\binom{N-K}{n-k}z^k$

$G'(z)=\dfrac1{\binom{N}{n}}\displaystyle\sum_{k=k_{\min}}^{k_{\max}} k\binom{K}{k}\binom{N-K}{n-k}z^{k-1}$

$G^{(k)}(z)=\dfrac1{\binom{N}{n}}\displaystyle\sum_{j=k_{\min}}^{k_{\max}} j^{\underline{k}}\binom{K}{j}\binom{N-K}{n-j}z^{j-k}$

